\chapter*{Preface}
\addcontentsline{toc}{chapter}{Preface}
\markboth{Preface}{Preface}

A genome is a stored description that no single cell reads in full. Different cells
carrying almost the same DNA express different parts of it, at different times and in
response to different histories. Populations change the description itself over
generations. Computing has its own versions of these separations: a program and its
running state, a model's weights and its activations, a catalogue of modules and the
plan selected for one request. This book asks whether making those separations explicit
and typed, the way genomes and cells appear to separate them, can give computational
organizations that are worth building. The question is \emph{whether}. The book does not
start from the assumption that they are.

\section*{Three meanings of genomic computation}
The phrase ``genomic computation'' is used for three different subjects, and much
confusion comes from sliding between them.

\begin{description}
\item[Computation about DNA.] Models that read recorded DNA sequence and predict
something about it, such as genomic language models, variant scorers and assay
predictors. Here DNA is \emph{data}.
\item[Computation with DNA.] Molecules as the physical substrate of computation:
strand-displacement circuits, molecular neural networks, algorithmic self-assembly and
DNA data storage. Here DNA is the \emph{machine}. The companion volume
\emph{DNA Computing: Molecular mechanisms, algorithms and evidence} develops this subject
from first principles~\cite{dnacomputingbook}.
\item[Computation inspired by DNA.] Digital architectures that borrow the
\emph{organization} of genomic systems: stored versus expressed structure, regulated
selection, strand symmetry, development and inheritance. Here DNA is a \emph{design
hypothesis}.
\end{description}

Evolutor works in the first and third senses and keeps both apart from the second. A
tokenized DNA string inside a neural network is not a molecule, and a software gate is not
a promoter. The Epilogue returns to all three senses and asks where they might
productively meet.

\section*{Three research directions with fixed names}
Throughout the book, \textbf{DOGMA} names a proposed non-Transformer, DNA-native research
direction together with its state-native engine; \textbf{Hermon DNA} names a
Transformer-based DNA research direction together with its attention-native engine; and
\textbf{Evolutor} names the broader theory, compiler, runtime, memory and governance layer
above both. The names are commitments about \emph{architecture}, not claims of success.
Each direction is reduced to a small executable reference model whose behaviour can be
checked exactly. The book then asks what experiment would show that the reference is worth
scaling, and what result would show that it is not.

\section*{What you will build}
Part~II builds the machinery of modern sequence learning from first principles:
objectives and legal information access, hand-derived gradients, reproducible training,
tokenization, recurrence, state-space scans, attention, the Transformer block and
conditional routing. Part~III maps real regulatory mechanisms to typed software
abstractions and stops each mapping where the biology ends. Parts~IV and~V specify DOGMA
and Hermon DNA. Part~VI fixes how unlike architectures can be compared without false
equivalence. Parts~VII and~VIII turn both references into serving engines (state pools,
checkpoints, batching, paged KV caches, continuous batching and speculative decoding)
while preserving exactly what each model computes. Part~IX builds Evolutor's runtime,
memory and compiler above both engines. Part~X covers deployment, isolation, distributed
placement and measurement. Part~XI turns application and capability claims, including
claims about general intelligence, into experimental designs. Part~XII collects
failures, open problems and the evidence record.

Almost every mechanism is implemented in a small reference program that you can run, and
almost every implementation is checked against an independent oracle. Wherever a
plausible bug exists, a deliberately broken negative control shows that the test can
actually catch it.

\section*{An evidence discipline}
The book separates kinds of evidence that are easy to blur: a theorem under stated
premises; a finite test of an implementation; a synthetic trace; a numerical emulation; a
task-trained result on held-out data; a measurement on real hardware; a biological
observation; and a hypothesis. Appendix~C summarizes what each can and cannot establish.
The same discipline applies to the book's own results. Several small trained pilots in
Parts~IV and~V stayed at or near chance under their declared budgets, and they are
reported as negative results rather than hidden. A complete manuscript, a passing test suite or an
elegant architecture diagram is not evidence of genomic usefulness, production readiness
or general intelligence. These distinctions do not dampen ambition. They are what will
make any eventual positive result credible.

\section*{Who this book is for}
The book is written for advanced undergraduates, graduate students, research engineers
and scientists who want to build sequence-model systems rather than only use them. It
assumes programming in Python, linear algebra, multivariable calculus and basic
probability. Prior deep-learning experience helps but is not required: Part~II derives
what later chapters need. No laboratory biology is assumed. Appendix~A gives the
molecular biology, coordinate conventions and data formats used in the genomic chapters.
The book contains no wet-laboratory procedures and no clinical interpretation of
variants.

\section*{Acknowledgments}
The general implementation-first progression of predicting, deriving, implementing,
testing and then examining consequences was informed by the author companion overview and
chapter map for Sebastian Raschka's \emph{Build a Large Language Model (From
Scratch)}~\cite{raschkaoverview}. The companion overview was inspected, not the full
commercial text. The prose, reference implementations, numerical examples, exercises and
illustrations of this book are its own unless explicitly attributed. Primary sources are
cited where their ideas are used, and each chapter's source notes record how deeply each
source was read.
